\paragraph{Kataoka's row $\lambda=1$.}
Here the three steps are $(2,2),(1,0),(0,1)$. Put
$X=\zeta^a$, $Y=\zeta^b$, $Z=XY$, and write $\nu=e\,v_2$ at conductor
$2^h$, where $e=2^{h-1}$. The Laplacian eigenvalue has the exact factorization
\[
 F=4-(Z+1)H,\qquad
 H=(X+Y)/Z + (Z-1)^2(Z+1)/Z^2 .
\]
Indeed the third pair of translated terms adds $-(Z-Z^{-1})^2$ to the
square-torus eigenvalue. Work first with $h\ge4$.
For mixed parities the square-torus eigenvalue has order $2$ and
$\nu((Z-Z^{-1})^2)=4$. The $2e$ mixed orbits therefore contribute $4e$.

For odd $a,b$ take $a=1$. If $b\equiv-1\pmod4$, the first term of
$H$ has order $2$, since $\nu(X+Y)=\nu(1-\zeta^{b-1+e})=2$.
The second has order at least $12$ or vanishes. Thus $\nu(H)=2$.
Over these $e/2$ orbits, the exponents $1+b+e$ range over all multiples
of four modulo $2^h$. For a nonzero such exponent, say $k$, the factor
$Z+1$ has order $2^{v_2(k)}\le e$. The finite orders sum to
$\sum_{j=2}^{h-1}(e/2^j)2^j=(h-2)e$. One orbit has $Z=-1$, giving
$\nu(F)=2e$. On the others $\nu(F)=\nu(Z+1)+2<2e$. These orbits
together contribute $(h+1)e-2$.

For $b\equiv1\pmod4$ both $Z-1$ and $Z+1$ have order $2$.
Consequently the second term of $H$ has order $6$.
The first has order $4$ for $b\equiv5\pmod8$ and at least $8$
(including infinity) for $b\equiv1\pmod8$, since $8\mid e$.
The two contributions to $H$ have distinct orders in both cases.
The resulting orders of $F$ are $6$ and $8$, each occurring $e/4$
times. Their sum is $7e/2$. We have proved
\[
          W_h=(h+17/2)2^{h-1}-2\qquad(h\ge4).
\]
Direct cyclotomic reduction at $h=1,2,3$ gives the respective shell sums
$7,16,48$. Thus $T_1=5$, $T_2=19$, $T_3=65$. The recurrence
$T_h-T_{h-1}=W_h-2$ now gives
\[
             T_n=(n+15/2)2^n-4n-7\qquad(n\ge3).
\]
